\section{Introduction}
\label{sec:introduction}

The Standard Model and general relativity are extraordinarily successful at organizing experimental facts.
At the same time, their combined fundamental inventory---multiple fields, symmetry sectors, symmetry breaking
structure, and many parameters fixed by measurement---invites a complementary line of inquiry:
whether a smaller set of ground ingredients could underlie this effective complexity.
This paper explores such a \emph{minimal-ingredient reconstruction} program.

\paragraph{Substrate-first strategy.}
We adopt a deterministic \emph{substrate-first} strategy: posit an ontically real wave medium and treat familiar
field theories as coarse-grained descriptions of constrained wave dynamics.
The goal is \emph{not} to dispute established phenomenology, but to ask what could make gauge structure,
relativistic kinematics, and particle-like quantization appear naturally as emergent regularities.

\paragraph{Ground ingredients (minimal commitments).}
The theoretical framework developed here keeps its primitive commitments deliberately small:
(i) a single substrate ontologically realised as a 4D cubic lattice embedded with codimension zero in a 4D Euclidean ambient, with the timelike direction picked out by the Lorentzian sign structure of the brane action (\Cref{subsec:ontology-kinematics}), and macroscopically described by its continuum long-wavelength theory on a spacelike slice,
(ii) pure geometric coupling between displacements perpendicular to a link and the lengthening of that link,
and (iii) a controlled separation between a linear wave sector and a nonlinear self-guidance sector.
Gauge potentials, relativistic coordinates, and particle labels are then interpreted as emergent bookkeeping attached to
stable excitation sectors rather than as independent fundamental substances.

\paragraph{Geometric source of nonlinearity.}
One modeling distinction is crucial throughout this paper.
The microscopic link law may be linear in the scalar extension, but the extension itself is the full Euclidean distance
between neighboring nodes in $\R^4$.
Equivalently, in the continuum description the induced metric already contains the quadratic term
$\partial_i u\,\partial_j u$ when the fourth embedding component is excited.
The primary nonlinearity of the theory therefore comes from geometry---the Pythagorean dependence of link lengths and the
induced metric---not from an added interaction field, an imposed threshold, or a hand-crafted saturation rule.
This point matters because the same geometric mechanism is intended to underwrite contraction, self-guidance, and
localized particle structure.

\paragraph{Substrate hypothesis.}
We assume a brane lattice with the 4D-in-4D ontology of \Cref{subsec:ontology-kinematics}: a 4D cubic lattice in a symmetric 4D Euclidean ambient, with the timelike direction selected by the brane action's Lorentzian sign structure.
For technical analysis on a spacelike slice the brane is parameterised by an embedding map
\begin{equation}
  \vecX:\Omega\times\R \to \R^4, \qquad (x,t)\mapsto \vecX(x,t), \qquad \Omega\subset\R^3,
  \label{eq:embedding-map}
\end{equation}
with $t$ the slicing parameter along the inside observer's timelike direction.
The dynamics is specified by a local hyperelastic action on $\Omega$, constructed from invariants of the
induced metric of the embedding.
This is the minimal coordinate-free way to write an elastic energy without hard-coding preferred directions into the
continuum description; microscopically, a cubic lattice substrate is allowed, and large-scale isotropy is treated as an
empirical constraint on admissible lattice couplings and parameter regimes \citep{Ogden1984,Holzapfel2000,DoCarmo1992}.

\paragraph{Gravity from transverse excitation (contraction).}
A central motivating idea is that excitation in the embedding dimension generically feeds back into
\emph{lateral} stresses.
If a region of the brane develops transverse structure $X^4=u(x,t)$, then the induced metric gains the quadratic
contribution $\partial_i u\,\partial_j u$.
Because the in-brane displacement is itself dynamical, a prestretched medium relaxes by generating an inward-biased
lateral response rather than only a local stress increase.
Coarse-grained over scales large compared to the carrier wavelength, such slowly varying contraction fields can bias the
propagation of other wave packets and the equilibrium of localized modes.
We therefore treat Newtonian gravity as an emergent, long-range \emph{contraction channel} mediated by transverse
excitation, not as a separately postulated fundamental field.

\paragraph{Per-wavepacket band isolation and geometric phase.}
The complex-envelope description that promotes the real lateral triplet to $\Psi\in\C^3$ does not require a
globally coherent carrier sector spanning macroscopic scales.
Instead it is applied per band-isolated excitation: each wavepacket carries its own local carrier frequency
$\omega_0$ and its own retained polarization subspace, defined as an eigenbundle of the linearized carrier
operator about the slowly varying background at the wavepacket's location.
Adiabatic transport of that local subspace induces a Berry (rank-1) or Wilczek--Zee (non-Abelian) connection,
which is interpreted as the natural origin of gauge variables in the per-wavepacket effective description
\citep{Pancharatnam1956,Berry1984,WilczekZee1984,Bliokh2015SpinOrbit,Cisowski2022}.

\paragraph{Microscopic cubicity and macroscopic isotropy.}
The cubic lattice hypothesis is not rotationally symmetric on the microscopic level.
The central question is whether long-wavelength observers built from the same substrate could nonetheless inherit a
universal effective metric --- and whether birefringence or residual direction dependence survives as a decisive diagnostic
if they do not.
The full dual-observer argument (lab anisotropy vs internal isotropy) is stated canonically in \Cref{para:dual-observer}.

\paragraph{Particles as solitons; probability not fundamental.}
We do not take the probabilistic interpretation of quantum theory as fundamental.
Instead, apparent point-like interactions and discrete outcomes are attributed to deterministic wave dynamics
combined with Fourier-localization constraints and nonlinear exchange events.
Particles are not fundamental point objects in this framework; they are localized standing-wave patterns of the
substrate, potentially with internal polarization and topological structure.
The spin-$\tfrac12$ $4\pi$ aspect is attributed to geometric phase/holonomy of an internal polarization
subspace rather than to literal multi-winding of a photon path.

\paragraph{Relation to existing ideas.}
The emergence of an effective light cone and Lorentz-like kinematics from a substrate whose time direction is selected by the action signature rather than by ambient geometry has close analogues in analogue-gravity systems \citep{Barcelo2011AnalogueGravity,Volovik2003HeliumDroplet,Jacobson2004EinsteinAether}.
The present use of the term ``brane'' is literal (an elastic medium), and should not be confused with
field-theoretic brane-world scenarios in high-energy theory \citep{RandallSundrum1999}.
The closest contemporary programme is the quaternionic quantum mechanics of Danielewski, Sapa, and Roth, in which
the vacuum is an ideal elastic solid (a Planck--Kleinert crystal), elementary particles are soliton-like standing
waves, gravity appears as a density-dependent change in wave speed (an effective index of refraction), and baryons
are modelled as three-quark standing-wave complexes \citep{Roth2023QQM2,DanielewskiRoth2024Baryons}.
We share the elastic-medium ontology and the particles-as-solitons commitment, but differ in three structural choices:
the carrier here is a \emph{real} vector substrate on a $4$D codimension-zero brane whose time direction is selected by
the action signature rather than an external time; the complex and gauge structure is \emph{derived} per wavepacket
from band isolation and Berry/Wilczek--Zee holonomy rather than postulated through a quaternion field; and we make no
claim to fix fundamental constants from the model.

\paragraph{Paper organization.}
\Cref{sec:scope} states the precise scope and the numerical priorities.
\Cref{sec:continuum-model} formulates the lattice and continuum substrate model, making the geometric source of the
nonlinearity explicit.
\Cref{sec:emergent-relativity} discusses the linear-wave regime and an emergent effective light cone.
\Cref{sec:geophase,sec:effective-field} develop the per-wavepacket band-isolation framing and the associated Berry/Wilczek--Zee gauge structure.
\Cref{sec:localized} formulates localized particle modes as self-adjoint eigenproblems with nonlinear self-guidance.
\Cref{sec:bell-interpretation} states the interpretive commitment of the theory under Bell's theorem: a retrocausal worldtube reading in which causality emerges from soliton chirality and entanglement is the continuity of one branching 4D worldtube.
\Cref{sec:discussion} then identifies the most decisive next step: a baryon-like triplet simulation on the cubic
substrate, where the axis-aligned internal structure is expected to matter most directly.
